% ==========================================================================
%  Chapter 83 -- Normative multiscale publication audit and redesign
% ==========================================================================

\chapter{Normative Multiscale Publication Audit and Redesign}
\label{ch:multiscale-publication-audit}

\section{Purpose and Normative Scope}
\label{sec:ms-pub-purpose}

This chapter is the normative reference for the public multiscale API in
\texttt{fall\_n}.  It supersedes the terminology and behavioural claims of
Chapters~\ref{ch:two-way-fe2} and \ref{ch:cyclic-validation} whenever those
chapters describe an earlier or partially implemented coupling loop.

The goal of the redesign is not merely to rename components, but to align the
software architecture with three simultaneous requirements:
\begin{enumerate}
  \item \textbf{physical coherence:} the public semantics must reflect the
    actual multiscale algorithm implemented in code;
  \item \textbf{architectural extensibility:} the hot path must remain
    zero-overhead while extension points stay explicit and SOLID at the API
    boundary;
  \item \textbf{publication readiness:} code, examples, tests, and
    documentation must tell the same story.
\end{enumerate}

\section{Official Nomenclature}
\label{sec:ms-pub-nomenclature}

The multiscale subsystem now exposes three officially distinct coupling modes:
\begin{description}
  \item[\texttt{OneWayDownscaling}] macro solve followed by micro solve, with no
    feedback to the macro model;
  \item[\texttt{LaggedFeedbackCoupling}] macro solve followed by micro solve and
    delayed feedback injection for subsequent macro steps;
  \item[\texttt{IteratedTwoWayFE2}] fixed-point macro--micro iteration with
    checkpoint/restore and a single accepted commit at the end of the step.
\end{description}

This separation is implemented in
\texttt{src/analysis/CouplingStrategy.hh} and reported explicitly through
\texttt{CouplingMode} in \texttt{src/analysis/MultiscaleTypes.hh}.  The
previous labels \texttt{OneWayCoupling} and \texttt{TwoWayStaggered} are
retained only as historical names in older documentation chapters.

\section{Audit of Prior Claims}
\label{sec:ms-pub-audit}

The redesign began by treating the existing code and documentation as
hypotheses rather than ground truth.  Table~\ref{tab:ms-defect-ledger}
records the main findings that motivated the current implementation.

\begin{table}[htb]
\centering
\small
\begin{tabular}{p{0.30\textwidth}p{0.14\textwidth}p{0.42\textwidth}}
\toprule
Claim or behaviour & Status & Normative conclusion \\
\midrule
The pre-redesign staggered driver was a full FE\textsuperscript{2} loop
  & incorrect
  & It was a lagged outer-loop driver: the macro step was solved outside the
    coupling iteration and was not restored/re-solved inside each fixed-point
    pass. \\
Tangents and forces were homogenised by the same operator
  & incorrect
  & The tangent came from boundary reactions while forces could be injected via
    a separate approximate path; the public API now packages both in a single
    \texttt{SectionHomogenizedResponse}. \\
Homogenised injection was local to the critical macro site
  & approximate
  & Beam elements supported blanket overrides on all section points; the public
    path is now \emph{site-based} through \texttt{CouplingSite}. \\
The driver owned no multiscale protocol details
  & incorrect
  & Older examples had to orchestrate \texttt{end\_of\_step()}, feedback
    injection, and convergence externally.  The orchestrator now owns the full
    lifecycle. \\
Regularisation was a physical constitutive choice
  & incorrect
  & The previous diagonal clamp was an implicit numerical stabiliser; it is now
    surfaced as an explicit regularisation policy with diagnostics. \\
All documentation under Chapters~74--82 described the current code
  & historical
  & Those chapters remain useful as design records, but this chapter is the
    authoritative statement of the current implementation. \\
\bottomrule
\end{tabular}
\caption{Defect ledger for the multiscale redesign.}
\label{tab:ms-defect-ledger}
\end{table}

\section{Underlying Mathematics and Physical Contract}
\label{sec:ms-pub-math}

\subsection{Macro generalised variables}

At the beam level, the section deformation vector is
\begin{equation}
  \mathbf{e}
  =
  \{\varepsilon_0,\kappa_y,\kappa_z,\gamma_y,\gamma_z,\theta'\}^{T},
\end{equation}
and the section resultant vector is
\begin{equation}
  \mathbf{s}
  =
  \{N,M_y,M_z,V_y,V_z,T\}^{T}.
\end{equation}

The local continuum sub-model is driven by Timoshenko-consistent end
kinematics extracted from the macro beam.  The current implementation uses the
kinematic transfer already established in Chapters~\ref{ch:multiscale_seismic}
and \ref{ch:nl-sub-model-evolver}, but the public contract is no longer
``beam element to entire sub-model'' in the abstract; it is
\emph{macro coupling site to local model} through \texttt{CouplingSite}.

\subsection{Normative homogenised response}

The public response object is
\texttt{SectionHomogenizedResponse}, containing:
\begin{equation}
  \mathbf{s}_0, \qquad
  \mathbf{D}_{\mathrm{hom}} = \frac{\partial \mathbf{s}}{\partial \mathbf{e}},
  \qquad
  \mathbf{e}_0.
\end{equation}
The macro-scale injection contract is
\begin{equation}
  \mathbf{s}(\mathbf{e})
  =
  \mathbf{s}_0 + \mathbf{D}_{\mathrm{hom}}(\mathbf{e}-\mathbf{e}_0),
  \label{eq:ms-pub-linearized-response}
\end{equation}
with \(\mathbf{e}_0\) stored as \texttt{strain\_ref}.  This eliminates the
ambiguity of ``tangent by one operator, forces by another'' in the public API.

\subsection{Boundary-reaction homogenisation}

The default production path in \texttt{NonlinearSubModelEvolver} now computes
section forces from assembled boundary reactions at the loaded end.  In compact
form,
\begin{equation}
  \mathbf{s}_{\mathrm{hom}}
  =
  \mathcal{H}_{\mathrm{br}}(\mathbf{u}_{\mu}),
\end{equation}
where \(\mathcal{H}_{\mathrm{br}}\) denotes the boundary-reaction operator and
\(\mathbf{u}_{\mu}\) is the converged micro displacement field.  The tangent is
obtained by perturbing the generalised deformation components and re-evaluating
the same reaction-based force operator:
\begin{equation}
  \mathbf{D}_{\mathrm{hom}}^{(:,j)}
  \approx
  \frac{\mathcal{H}_{\mathrm{br}}(\mathbf{u}_{\mu}^{(j,+)})
        - \mathcal{H}_{\mathrm{br}}(\mathbf{u}_{\mu})}{h_j}.
\end{equation}

Volume averaging is retained only as a comparison/validation route and is
reported through \texttt{HomogenizationOperator::VolumeAverage}; it is not the
default FE\textsuperscript{2} production operator.

\subsection{Convergence criterion}

The default coupling convergence is no longer defined solely through a tangent
matrix norm.  The implemented criterion monitors:
\begin{equation}
  r_f^{(i)}
  =
  \frac{\lVert \mathbf{s}_{\text{macro}}^{(i)}
  - \mathbf{s}_{\text{micro}}^{(i)} \rVert}
  {\max\left(1,\lVert \mathbf{s}_{\text{macro}}^{(i)} \rVert,
                \lVert \mathbf{s}_{\text{micro}}^{(i)} \rVert\right)},
\end{equation}
for each coupled site \(i\), together with an analogous relative tangent
residual.  The default policy in
\texttt{src/analysis/CouplingStrategy.hh} is
\texttt{ForceAndTangentConvergence}.

\section{Implemented Architecture}
\label{sec:ms-pub-architecture}

\subsection{Public data contracts}

The redesign introduces explicit public types:
\begin{itemize}
  \item \texttt{CouplingSite}: identifies the exact macro site
    (\texttt{macro\_element\_id}, section Gauss point, \(\xi\), local frame);
  \item \texttt{MacroSectionState}: local macro strain/resultant snapshot;
  \item \texttt{SectionHomogenizedResponse}: homogenised forces, tangent,
    reference strain, operator provenance, regularisation diagnostics,
    perturbation sizes, tangent-column validity, and finite-difference
    provenance;
  \item \texttt{CouplingIterationReport}: mode, termination reason,
    iteration count, residuals, hard-failure sites, rollback state,
    regularisation presence, and macro/micro timing;
  \item \texttt{ModelCheckpoint} and solver-specific
    \texttt{checkpoint\_type}: exact rollback state for FE\textsuperscript{2}
    iteration.
\end{itemize}

The typed path through \texttt{MultiscaleModel<MacroBridgeT, LocalModelT>} is
the normative public surface.  The type-erased
\texttt{LocalModelHandle} remains available only as an experimental utility
until erased checkpoint/restore semantics are completed.  Likewise,
\texttt{src/reconstruction/HomogenisationStrategy.hh} is retained as an
internal validation utility rather than a recommended public extension point.

\subsection{SOLID without sacrificing the hot path}

The multiscale orchestration now follows the design principle ``static in the
kernel, explicit at the boundary'':
\begin{itemize}
  \item \textbf{Single Responsibility Principle.}
    \texttt{BeamMacroBridge} is responsible for macro extraction/injection;
    \texttt{MultiscaleModel} aggregates coupling sites and local models;
    \texttt{MultiscaleAnalysis} owns the protocol; the local solver owns local
    constitutive evolution.  Internally, the nonlinear sub-model path now
    isolates boundary-condition application, persistent state lifecycle, and
    boundary-reaction homogenisation into dedicated header-only components.
  \item \textbf{Open/Closed Principle.}
    Coupling algorithms, convergence criteria, relaxation policies, and
    executors are all pluggable through small interfaces.
  \item \textbf{Liskov Substitution Principle.}
    Macro solvers participate through the
    \texttt{CheckpointableSteppableSolver} concept rather than through ad hoc
    assumptions about \texttt{DynamicAnalysis} or \texttt{NonlinearAnalysis}.
  \item \textbf{Interface Segregation Principle.}
    The local-model concept requires only what the multiscale loop actually
    needs: kinematics update, solve, homogenised response, commit/revert, and
    checkpoint control.
  \item \textbf{Dependency Inversion Principle.}
    \texttt{MultiscaleAnalysis} depends on abstract coupling policies and on
    concepts, not on concrete drivers.
\end{itemize}

These abstractions are intentionally placed at the orchestration layer.  The
inner loop remains template-based and avoids virtual dispatch in the hot path.
This preserves the original performance philosophy of \texttt{fall\_n} while
making the multiscale subsystem externally reusable.

At the local-solver level, the same principle is now applied internally:
\begin{itemize}
  \item \texttt{LocalBoundaryConditionApplicator} owns beam-to-continuum
    kinematic updates and writes the imposed boundary vector;
  \item \texttt{PersistentLocalStateOps} owns synchronisation,
    commit, and revert of the persistent continuum state;
  \item \texttt{BoundaryReactionHomogenizer} owns the section-force extraction,
    tangent diagnostics, and response classification.
\end{itemize}
These components are value-type, header-only helpers with pointer-based
composition; they add no virtual dispatch to the micro hot path while creating
an explicit seam for future replacement of the finite-difference tangent by a
more rigorous consistent linearisation.

\subsection{Macro and local rollback}

The FE\textsuperscript{2} loop is only meaningful if the macro and micro
states can both be restored exactly.  The current implementation provides:
\begin{itemize}
  \item model-level snapshots through \texttt{ModelCheckpoint};
  \item solver-level snapshots in \texttt{DynamicAnalysis} and
    \texttt{NonlinearAnalysis};
  \item local-model snapshots in \texttt{NonlinearSubModelEvolver};
  \item explicit trial/commit control via \texttt{set\_auto\_commit(false)}
    during iteration and a single \texttt{commit\_trial\_state()} after
    acceptance.
\end{itemize}

\section{Implemented Iterated Two-Way FE\textsuperscript{2} Algorithm}
\label{sec:ms-pub-algorithm}

Algorithm~\ref{alg:iterated-two-way-fe2} summarises the implemented
orchestrator semantics.

\begin{algorithm}[htb]
\DontPrintSemicolon
\caption{Implemented \texttt{IteratedTwoWayFE2} protocol}
\label{alg:iterated-two-way-fe2}
\KwIn{macro solver \(S_M\), multiscale model \(\mathcal{M}\),
      predictor response \(\mathbf{R}^{(0)}\)}
\KwOut{accepted macro step and committed local states}
\BlankLine
capture macro checkpoint \(C_M\)\;
capture one checkpoint \(C_\mu^{(i)}\) for each local model\;
disable automatic macro and local commit\;
\For{\(k = 0,1,\dots,k_{\max}-1\)}{
  restore \(C_M\) and all \(C_\mu^{(i)}\)\;
  inject predictor response at each \texttt{CouplingSite}\;
  advance one macro step\;
  extract local beam kinematics for each site\;
  solve local models in parallel\;
  build \texttt{SectionHomogenizedResponse} for each site\;
  compute force and tangent residuals\;
  apply relaxation policy if configured\;
  \If{\(k \ge 1\) and all convergence criteria are satisfied}{
    accept current iteration\;
    \textbf{break}\;
  }
  set predictor \(\leftarrow\) current response\;
}
inject accepted response at each site\;
commit macro trial state\;
commit local trial states and call \texttt{end\_of\_step()} once\;
\end{algorithm}

Two implementation decisions are deliberate:
\begin{enumerate}
  \item the first iteration is always treated as a predictor pass and is never
    accepted immediately;
  \item \texttt{end\_of\_step()} is owned by the orchestrator, not by the
    driver.
\end{enumerate}

\section{HPC and Build Surface}
\label{sec:ms-pub-hpc}

The first publication-ready HPC layer focuses on robust shared-memory
parallelism:
\begin{itemize}
  \item \texttt{SerialExecutor} provides the reference path;
  \item \texttt{OpenMPExecutor} parallelises local-model solves over coupling
    sites;
  \item \texttt{CouplingCommunicators} makes the future MPI ownership explicit
    without pretending that a distributed FE\textsuperscript{2} engine is
    already complete.
\end{itemize}

The build system now exposes a dedicated public CMake target:
\begin{center}
\texttt{fall\_n::multiscale\_api}
\end{center}
and a slim public umbrella header:
\begin{center}
\texttt{src/analysis/MultiscaleAPI.hh}
\end{center}
which gathers the multiscale-specific headers without routing users through the
project-wide \texttt{header\_files.hh} umbrella.

\section{Verification and Test Matrix}
\label{sec:ms-pub-tests}

The current redesign is supported by the following regression and contract
tests:
\begin{itemize}
  \item \texttt{tests/test\_multiscale\_analysis\_api.cpp}:
    semantic separation of coupling modes, ownership of local lifecycle,
    analysis-local coupling start, accepted FE\textsuperscript{2} commit,
    explicit rollback semantics, hard-failure classification, and parity
    between serial and OpenMP execution on the API harness;
  \item \texttt{tests/test\_micro\_solve\_executor.cpp}:
    serial/OpenMP executor contract and communicator defaults;
  \item \texttt{tests/test\_beam.cpp}:
    section-level beam behaviour, including the element path touched by local
    homogenised injection;
  \item \texttt{tests/test\_steppable\_solver.cpp} and
    \texttt{tests/test\_steppable\_dynamic.cpp}:
    solver contracts relied upon by the orchestrator;
  \item \texttt{tests/test\_evolver\_advanced.cpp}:
    persistent local-model evolution, adaptive finite-difference tangent
    diagnostics, pure axial/flexural/shear/torsional regression cases,
    linearised consistency, reaction-versus-volume operator comparison, and
    axial power consistency.
\end{itemize}

In addition, the heavy executable targets
\texttt{fall\_n\_lshaped\_multiscale},
\texttt{fall\_n\_lshaped\_multiscale\_16},
\texttt{fall\_n\_table\_multiscale}, and
\texttt{fall\_n\_table\_cyclic\_validation}
serve as integration-level build checks of the public multiscale surface.

\section{Known Limitations}
\label{sec:ms-pub-limitations}

The redesign closes the most serious semantic gaps, but several publication
limitations remain explicit:
\begin{itemize}
  \item the default tangent is now obtained by adaptive finite differences with
    a central-difference preference and one-sided fallback, not by a fully
    consistent algorithmic linearisation;
  \item the accepted FE\textsuperscript{2} step is a fixed-point iteration, not
    a monolithic multilevel Newton solve;
  \item MPI-distributed local-model execution is not yet implemented beyond the
    communicator contract;
  \item \texttt{NonlinearSubModelEvolver} still contains several responsibilities
    internally, but the boundary-condition, state-lifecycle, and
    homogenisation concerns have already been split from the main class body;
  \item the GitHub Actions workflow added for multiscale stability has been
    authored and documented, but only the local PowerShell gate has been
    executed in this audit cycle;
  \item historical chapters still contain earlier design narratives and should
    be read as provenance, not as the normative behavioural specification.
\end{itemize}

\section{Future Work}
\label{sec:ms-pub-future}

The next publication-strength milestones are:
\begin{enumerate}
  \item replace the adaptive finite-difference tangent by a fully consistent
    algorithmic tangent while retaining the same public response contract;
  \item finish the internal split of \texttt{NonlinearSubModelEvolver} from
    the current state into a solver core, diagnostics service, and output
    service, reusing the already extracted boundary-condition, state, and
    homogenisation helpers;
  \item add energy-consistency benchmarks and distributed MPI execution of
    local models;
  \item package installable public headers for the multiscale target instead of
    exposing the source-tree include path;
  \item extend the coupling-site abstraction beyond beam sections to shell or
    continuum macro sites.
\end{enumerate}

\section{Conclusion}
\label{sec:ms-pub-conclusion}

The multiscale subsystem of \texttt{fall\_n} is now documented and exposed as a
library feature rather than as a driver convention.  The redesign does not
claim that every research-grade FE\textsuperscript{2} feature is complete; it
does claim, and the current code supports, that:
\begin{enumerate}
  \item the coupling modes are semantically distinct and named honestly;
  \item iterated two-way FE\textsuperscript{2} performs real macro/micro
    checkpointed iteration;
  \item forces, tangents, and reference strain are injected through one public
    response object;
  \item the lifecycle no longer leaks into the examples;
  \item the documentation now states these facts explicitly.
\end{enumerate}
